%%%%%%%%%%%%%%%%%%%%%%%%%%%%%%%%%%%%%%%%%%%%%%%%%%%%%%%%%%%%%%
%%%%%%%%%%%%%%%%%% MA-0815 Análisis armónico %%%%%%%%%%%%%%%%%
%%%%%%%%%%%%%%%%%%%%%%%%%%%%%%%%%%%%%%%%%%%%%%%%%%%%%%%%%%%%%%
% Propuesta de programa adaptada de:
% Posgrado en Matemática, UCR, II-2019,
% «Propuesta de Programa: Análisis Armónico».
% Referencia principal de contenidos: Duoandikoetxea \cite{Duo01}.
\newpage
\subsection*{MA-0815 Análisis armónico}
\addcontentsline{toc}{subsection}{MA-0815 Análisis armónico}

\hypertarget{MA0815}{}
\begin{refsection}
\begin{table}[H]
\centering{
\begin{tabular}{|ll|}
\hline
\textbf{Año:} IV  & \textbf{Requisitos:} MA-0705 (o MA-0625) Análisis Real I \\
\textbf{Ciclo:} Optativa  & \textbf{Correquisitos:} Ninguno \\
\textbf{Tipo de curso:} Teórico & \textbf{Horas presenciales por semana:} 5 \\
\textbf{Créditos:} 5 & \textbf{Horas de trabajo independiente por semana:} 10 \\
\textbf{Virtualidad:} P  &  \\ \hline
\end{tabular}
}
\end{table}

\subsubsection*{Descripción del curso}

Este curso optativo introduce el análisis armónico de variable real: el estudio de la transformada de Fourier, de operadores maximales y de integrales singulares sobre $\R^n$. Se ubica en el cuarto año, una vez que la persona estudiante dispone de la medida e integración de Lebesgue y de una primera aproximación a los espacios $L^p$ (MA-0705 Análisis Real~I). Complementa la introducción a series y transformadas de Fourier que puede haberse visto en MA-0725 Análisis Real~II, y se articula de manera natural con el análisis funcional y con las ecuaciones diferenciales parciales.

El programa se articula a partir de la propuesta de Análisis Armónico del Posgrado en Matemática de la Universidad de Costa Rica (II-2019) y sigue, en lo esencial, el desarrollo de Duoandikoetxea \cite{Duo01}: espacios $L^p$ e interpolación, transformada de Fourier, funciones maximales, transformadas de Hilbert y de Riesz, integrales singulares de Calderón--Zygmund y teoría de Littlewood--Paley.

El énfasis es teórico: se requieren definiciones precisas, enunciados completos y demostraciones de los teoremas fundamentales. Los ejemplos en $L^p(\R^n)$, en el círculo y en el espacio de Schwartz se usan para concretar la teoría, no para sustituirla. Se espera que la persona estudiante formule argumentos rigurosos, construya ejemplos y contraejemplos, y reconozca el alcance de cada hipótesis. Haber cursado MA-0725 Análisis Real~II o MA-0806 Análisis funcional es recomendable, aunque no es requisito formal.

\subsubsection*{Objetivos}

\textbf{General:} Que la persona estudiante sea capaz de resolver problemas y desarrollar teoría a partir de los fundamentos del análisis armónico de variable real en $\R^n$.

\textbf{Específicos:} Al finalizar el curso se espera que la persona estudiante sea capaz de:

\begin{enumerate}
\item Emplear las propiedades de los espacios $L^p$ y $L^{p,\infty}$, y los teoremas de interpolación de Riesz--Thorin y de Marcinkiewicz, para acotar operadores.
\item Definir y manipular la transformada de Fourier en $L^1$, en la clase de Schwartz y en el sentido de las distribuciones temperadas, y extenderla a $L^p$ para $1<p\le 2$.
\item Analizar aproximaciones de la identidad y la función maximal de Hardy--Littlewood, y aplicar el teorema maximal a problemas de convergencia y diferenciación.
\item Construir las transformadas de Hilbert y de Riesz, relacionarlas con la función analítica y justificar su acotación en $L^p$.
\item Demostrar la acotación $L^2$ de integrales singulares homogéneas y la acotación débil $(1,1)$ mediante la descomposición de Calderón--Zygmund, y deducir la acotación en $L^p$.
\item Interpretar funciones cuadradas y multiplicadores de Fourier (Marcinkiewicz, Hörmander--Mihlin) en el marco de la teoría de Littlewood--Paley.
\item Consultar literatura especializada y comunicar argumentos de análisis armónico con rigor, tanto de forma oral como escrita.
\end{enumerate}

\subsubsection*{Contenidos}

\begin{enumerate}
\item \textbf{Espacios $L^p$ (3 semanas):}
\begin{enumerate}
    \item Preliminares: funcionales lineales, acotación y dualidad.
    \item Definición y propiedades de los espacios $L^p$.
    \item Dualidad de los espacios $L^p$.
    \item Espacios $L^p$-débiles.
    \item Interpolación de espacios $L^p$: teoremas de Riesz--Thorin y de Marcinkiewicz.
\end{enumerate}

\item \textbf{La transformada de Fourier (3 semanas):}
\begin{enumerate}
    \item Transformada de Fourier en $L^1$.
    \item Clase de Schwartz y distribuciones temperadas.
    \item Transformada de Fourier en $\mathcal{S}$ y en $\mathcal{S}'$.
    \item Transformada de Fourier en $L^p$, $1<p\le 2$.
    \item Convergencia y sumabilidad de integrales de Fourier.
\end{enumerate}

\item \textbf{Funciones maximales (2 semanas):}
\begin{enumerate}
    \item Aproximaciones de la identidad.
    \item Función maximal de Hardy--Littlewood.
    \item Función maximal fuerte.
    \item Función maximal diádica.
\end{enumerate}

\item \textbf{Las transformadas de Hilbert y de Riesz (2 semanas):}
\begin{enumerate}
    \item Construcción y propiedades.
    \item Relación con las funciones analíticas.
    \item Acotación en $L^p$. Teoremas de Riesz y de Kolmogorov.
    \item Transformadas de Riesz.
\end{enumerate}

\item \textbf{Integrales singulares de tipo convolución (3 semanas):}
\begin{enumerate}
    \item Integrales singulares homogéneas y maximales.
    \item Acotación en $L^2$.
    \item Método de rotaciones.
    \item Descomposición de Calderón--Zygmund.
    \item Operadores de Calderón--Zygmund.
    \item Acotación débil $(1,1)$ de integrales singulares.
    \item Condiciones para la acotación en $L^p$.
\end{enumerate}

\item \textbf{Teoría de Littlewood--Paley (2 semanas):}
\begin{enumerate}
    \item Desigualdades en valores vectoriales.
    \item Funciones cuadradas y $L^p$.
    \item Multiplicadores de Marcinkiewicz.
    \item Multiplicadores de Hörmander--Mihlin (\emph{ampliación}).
    \item Multiplicadores de Bochner--Riesz (\emph{ampliación}).
    \item Función maximal esférica (\emph{ampliación}).
\end{enumerate}

\end{enumerate}

\subsubsection*{Metodología}

\noindent \textbf{Lineamientos metodológicos:} La persona docente a cargo de este curso debe respetar las siguientes pautas:
\begin{enumerate}
    \item La presentación del material requiere de definiciones precisas, enunciados completos y demostraciones de los teoremas fundamentales (interpolación, teorema maximal, acotación de Hilbert--Riesz y descomposición de Calderón--Zygmund), aunque también es importante brindar contexto y motivación.
    \item Cada construcción abstracta debe acompañarse de al menos un ejemplo concreto en $L^p(\R^n)$, en el círculo o en la clase de Schwartz que ilustre el alcance de las hipótesis.
    \item Cada estudiante debe realizar una revisión de la bibliografía como complemento al trabajo presencial.
    \item Se fomenta la comunicación clara y rigurosa de argumentos, tanto oral como escrita.
\end{enumerate}

\textbf{Sugerencias metodológicas:} Adicionalmente, se tienen las siguientes recomendaciones para la persona docente a cargo de este curso:
\begin{enumerate}
    \item Tomar \cite{Duo01} como hilo conductor y recurrir a \cite{Gra14} o a \cite{SW71} cuando se desee mayor detalle en interpolación, integrales singulares o Littlewood--Paley.
    \item Explicitar el puente con el curso de análisis real: los espacios $L^p$ y la diferenciación de Lebesgue dejan de ser un apéndice de la integral y pasan a ser el lenguaje de los operadores.
    \item Dado que algunos temas permiten cubrirse con distintos grados de profundidad, se recomienda no sobrecargar al estudiantado. Entre estos temas están:
    \begin{itemize}
        \item Las pruebas completas de los teoremas de interpolación (pueden enunciarse y aplicarse, con una de las dos demostrada).
        \item El tratamiento sistemático de las distribuciones temperadas.
        \item Los multiplicadores de Hörmander--Mihlin y de Bochner--Riesz, y la función maximal esférica.
    \end{itemize}
    \item Promover el trabajo constante mediante listas de ejercicios que combinen demostraciones, cálculo de núcleos y construcción de ejemplos.
    \item Incorporar un proyecto o exposición breve sobre un tema de ampliación o sobre una aplicación a ecuaciones diferenciales o al análisis funcional.
    \item Sugerir el uso de \LaTeX{} para la presentación de tareas y del proyecto.
\end{enumerate}

\nocite{Duo01}
\nocite{Fol99}
\nocite{Gra14}
\nocite{Gra14b}
\nocite{Kat04}
\nocite{Ste93}
\nocite{SS05}
\nocite{SW71}

\printbibliography[heading=subbibliography,section=\the\value{refsection}]
\end{refsection}
